\documentclass{article}
\usepackage{tikz}
\usetikzlibrary{angles, quotes}
\usepackage{amssymb}
\usepackage{pgfplots}
\pgfplotsset{width=10cm,compat=1.9}
\usepgfplotslibrary{statistics}
\usepackage[left=2.54cm,right=2.54cm,top=2.54cm,bottom=2.54cm]{geometry}
\usepackage{tasks} %Paquete necesario para la producción de items horizontales para algunos ejercicios de tarea empleando el comando /begin{multicols}}{#}.../end{multicols} para ayudarnos a reducir las páginas
\usepackage{pdflscape} %comando que permite cambiar a orientación horizontal las hojas%
\usepackage{soul} %Paquete para permitir la compilación de acentos usando el comando {\hl{}}}\\
\sethlcolor{yellow(munsell)} %comando para definir el color para el subrayado usando el comando \hl
\usepackage[utf8]{inputenc}
\usepackage[spanish]{babel}
\usepackage{icomma}
\usepackage{ragged2e}		
\usepackage{siunitx}
\usepackage{url}
\usepackage[colorlinks=true, urlcolor=blue,  linkcolor=Cerulean, citecolor=green]{hyperref}
\usepackage{pdfpages}
\usepackage{blindtext} %Paquete que genera texto ficticio (dummy text) con el comando: \blindtext
\usepackage{cancel} %in the preamble gives you four different modes of striking through: \cancel{text to cancel} draws a diagonal line (slash) through its argument, \bcancel{text to cancel} uses the negative slope (a backslash), \xcancel{text to cancel} draws an X (actually \cancel plus \bcancel), \cancelto{〈value〉}{〈expression〉} draws a diagonal arrow through the 〈expression〉pointing to the 〈value〉 (math-mode only)
\usepackage{csquotes}
\usepackage{afterpage}
\usepackage{parskip} 
\usepackage{float}
\usepackage{enumitem}
%\usepackage{multicol}%Paquete que permite la creación aislada de columnas de texto. De esta forma se reduce la cantidad de páginas en nuestro documento
\usepackage{lipsum} %Paquete que genera texto de relleno al igual que \usepackage{blindtext}, se usa con el comando \lipsum[#-#] los #(hashtags) delimitan el número de páginas que deseamos ocupar del paquete lipsum para usarlos como relleno. 

% Definir el comando personalizado para la nota
\newcommand{\mynote}[1]{%
\begin{tikzpicture}[baseline=-0.75ex]
\node[draw, circle, fill=Apple Green!60, inner sep=2pt] (note) {\textbf{Nota:}};
\end{tikzpicture}%
\ \textit{#1}%
}

% Definición de comandos para variables con barra
\newcommand{\xbar}{\overline{x}}
\newcommand{\ybar}{\overline{y}}

\newenvironment{Figura}
{\par\medskip\noindent\minipage{\linewidth}}
{\endminipage\par\medskip}
\usepackage{caption}
\usepackage[
backend=biber,
sorting=none,
url=true
]{biblatex} %bibliografía
\addbibresource{biblio.bib}
% Establecer el espacio entre las entradas de la bibliografía
\setlength{\bibitemsep}{1\baselineskip} % Puedes ajustar el valor según tus preferencias
\usepackage{amsmath,amsthm,amssymb,amsfonts}
\usepackage{pifont} %Permite la compilación del comando \xmark: es el símbolo de la tachita
\usepackage{mathtools} %permite la compilación de símbolos para matrices
\usepackage{empheq} %Paquete que se relaciona con \usepackage[most]{tcolorbox} para la creación de cuadros/cajas de colores para delimitar los resultados matemáticos o líneas de texto usando la línea de comando: \begin{empheq}[box={\mymath[colback=orange(webcolor)!70,drop lifted shadow]}]{equation*} \end{empheq}
\usepackage[most]{tcolorbox} %Paquete que se relaciona con \usepackage{empheq} para la creación de cuadros/cajas de colores para delimitar los resultados matemáticos o líneas de texto
\renewcommand{\qedsymbol}{$\blacksquare$}
\usetikzlibrary{positioning,decorations.pathreplacing} %Paquete que permite la compilación de llaves para cuadros sinópticos (brace diagram) 
\usepackage{schemata} %The schemata package is designed just for these brace diagrams
\usepackage{booktabs}

\newcommand\AB[2]{\schema{\schemabox{#1}}{\schemabox{#2}}} %comando o paquete necesario para crear cuadros sinópticos

\newtcbox{\mymath}[1][]{%
nobeforeafter, math upper, tcbox raise base,
enhanced, colframe=blue!30!black,
colback=blue!30, boxrule=1pt,
#1} %Comando importante para encasillar los resultados matemáticos en cajas de diferentes colores, se relaciona con los paquetes \usepackage{empheq} y \usepackage[most]{tcolorbox} 

\newenvironment{sysmatrix}[1]
{\left(\begin{array}{@{}#1@{}}}
{\end{array}\right)}
\newcommand{\ro}[1]{%
\xrightarrow{\mathmakebox[\rowidth]{#1}}%
}
\newlength{\rowidth}% row operation width
\AtBeginDocument{\setlength{\rowidth}{3em}}  %Comando importante para laproducción de lineas de operaciones entre matrices, método de Gauss o Gauss Jordan se relaciona con \newenvironment{sysmatrix}

%formato para cambiar el horario a español
\usepackage[spanish]{datetime2}
\DTMsetdatestyle{spanish}

\renewcommand{\today}{\DTMdisplaydate{\the\year}{\the\month}{\the\day }{-1}}
%%%%%%%%%%%%%%%%%%%%%%%%%%%%%%%%%%%%%%%%%

\usepackage{datetime}
\newcommand{\mycurrenttime}{\xxivtime}

%%%%%%%%%%%%%%%%%%%%%%%%%%%%%%%%%%%%%%%%%%%%%%%%%%%%%%%%%%%%%%%%%%%%%%%%%%%%%%%%%%%%%%%%%%%%%%%%%%%%%%%%%%%%%%%%%%%%%%%%%%%%%%%%%%%%%%%%%%%%%%%%%


%%%%%%%%%%%%%Caja de color para el título%%%%%%%%%%%%%%%%%%%%%%%%%%%%%%%%%%%%%%%%%%%%%%%%%%%%%%

\definecolor{myframecolor}{RGB}{1, 45, 75} % Prussian Blue
\definecolor{myboxcolor}{RGB}{0, 107, 92} % Apple Green

\newtcolorbox{mybox}{
enhanced,
colback=myboxcolor!7, % Color de fondo del cuadro
colframe=myframecolor, % Color del marco del cuadro
arc=0pt,
boxrule=1pt,
borderline west={2mm}{-10mm}{myframecolor}, % Borde en el lado izquierdo
sharp corners=southwest,
width=\linewidth,
before=\par\vspace{\bigskipamount}, % Espacio antes del cuadro
after=\par\vspace{\bigskipamount} % Espacio después del cuadro
}

%%%%%%%%%%%%%%%%%%%%%%%%%%%%%%%%%%%%%%%%%%%%%%%%%%%%%%%%%Nueva caja%%%%%%%%%%%%%%%%%%%%%%%%%%%%%%%%%%%%%%%%%%%%%%%%%%%%%%%%%%%%%%%%%

\newtcolorbox{mybox2}{
enhanced,
colback=Apple Green!7, % Color de fondo del cuadro
colframe=Apple Green, % Color del marco del cuadro
arc=0pt,
boxrule=1pt,
borderline west={2mm}{-10mm}{Apple Green}, % Borde en el lado izquierdo
sharp corners=southwest,
width=\linewidth,
before=\par\vspace{\bigskipamount}, % Espacio antes del cuadro
after=\par\vspace{\bigskipamount} % Espacio después del cuadro
}


\newcommand{\euler}{e} %Comando para producir letra e de euler
\newenvironment{remark}{\par\vfill\footnotesize % Vertical white space above the remark and smaller font size
\begin{list}{}{
	\leftmargin=80pt % Indentation on the left
	\rightmargin=60pt}\item\ignorespaces % Indentation on the right
\makebox[-2.5pt]{\begin{tikzpicture}[overlay]
		\node[draw=Horizon!60,line width=2.5pt,circle,fill=Horizon!25,font=\sffamily\bfseries,inner sep=4pt,outer sep=0pt] at (-19pt,5pt){\textcolor{Horizon}{Nota}};\end{tikzpicture}} % Blue Nota in a circle
\advance\baselineskip -1pt}{\end{list}\vskip5pt} % Tighter line spacing and white space after remark
\usepackage{graphicx}

\usepackage{titling}

\input{colores.tex}

%%Producir signos de cita textual``Asignatura''

\newcommand{\xmark}{\ding{55}} %%%comando que genera la tachita

\newenvironment{MyColorPar}[1]{%
\leavevmode\color{#1}\ignorespaces%
}{%
}%

%%%%%%%%%%%%%%%%%%% Título de la tarea, Nombre de alumno

\title{Clase \# 43 - Estado Sólido}
\author{\emph{JA} $\boldsymbol{\mid}$}
\date{\today}

%%%%%%%%%%%%%%%%%%% Datos de la Materia

\usepackage{fancyhdr}
\fancypagestyle{plain}{%  the preset of fancyhdr 
\fancyhf{} % clear all header and footer fields
\fancyfoot[R]{
\includegraphics[width=2cm]{LogoFCFMBUAP (1).png}}
\fancyfoot[L]{{\bfseries{\thedate{}}} a las {\bfseries{\mycurrenttime{}}} horas{} (GMT-6, H. Puebla de Zaragoza, Pue)}
\fancyhead[L]{Estado sólido}
\fancyhead[R]{\theauthor}
}
\makeatletter
\def\@maketitle{%
\newpage
\null
\vskip 1em%
\begin{center}%
\let \footnote \thanks
{\LARGE \@title \par}%
\vskip 1em%
%{\large \@date}%
\end{center}%
\par
\vskip 1em}
\makeatother

\usepackage{lipsum}  
\usepackage{cmbright}

%%%%%%%%%%%%%%%%%%%%%%%%%%%%%%%%%%%%%%%%%%%%%%%%%%%%%%%%%%%%%%%%%%%%%%%%%%%%%%%%%%%%%%%%%%%%%%%%%%%%%%%%%%%%%%%%%%%%%%%%%%%%%%%%%%%%%%%%%%%%%%%%%%%%%%%%%%%%%%%%%%%%%%%%%%%%%%%%%%%%%%%%%%%%%%%%%%%%%%

% Por si el color se llama "Apple Green" con espacio en colores.tex:
\colorlet{AppleGreen}{Apple Green}

% Caja AppleGreen con tono ajustable
\newenvironment{AGbox}[1][7]% 7 = tono por defecto (más grande => más claro)
{%
  \begin{tcolorbox}[
    enhanced, breakable,
    colframe=AppleGreen,
    colback=AppleGreen!#1,
    boxrule=1pt,
    borderline west={2mm}{-10mm}{AppleGreen},
    sharp corners=southwest,
    left=14pt,right=14pt,top=10pt,bottom=10pt,
     boxsep=6pt,
  ]%
}{%
  \end{tcolorbox}%
}



%%%%%%%%%%%%%%%%%%%%%%%%%%%%%%%%%%%%%%%%%%%%%%%%%%%%%%%%%%%%%%%%%%%%%%%%%%%%%%%%%%%%%%%%%%%%%%%%%%%%%%%%%%%%%%%%%%%%%%%%%%%%%%%%%%%%%%%%%%%%%%%%%%%%%%%%%%%%%%%%%%%%%%%%%%%%%%%%%%%%%%%%%%%%%%%%%%%%%%

\begin{document}



\maketitle


%%%%%%%%%%%%%%%%%%% Datos del profesor, campus y aula

\begin{mybox}
\noindent\begin{tabular}{@{}ll}
	{\bfseries{Profesor}} & Dra. Ana Lilia González Ronquillo\\
	\textcolor{prussianblue}{\bfseries{Campus}} / \textcolor{trueblue}{\bfseries{Facultad}}  & \textcolor{prussianblue}{\bfseries{C.U. BUAP}} /  \textcolor{trueblue}{\bfseries{FCFM}} \\
	\textcolor{prussianblue}{\bfseries{Edificio}} /  \textcolor{trueblue}{\bfseries{Salón}}    & 	\textcolor{prussianblue}{\bfseries{FM9}} /  \textcolor{trueblue}{\bfseries{301}}   
\end{tabular} 
\end{mybox} \vspace{1cm}

\tableofcontents

\vspace{2cm}


\section{Clase de hoy: Electrones bajo un potencial periódico \underline{débil}}

Tenemos:

\[  - \frac{\hbar^{2}}{2m} \frac{d^{2}\psi}{dx^{2}} + V(x) \psi = \epsilon \psi  \]


Tenemos:

\[ \psi \to \psi_{\vec{k}}(x) \]

\[ k_{0}^{2} \equiv \frac{2m\epsilon}{\hbar^{2}}; \quad \text{y} \quad - \frac{2m}{\hbar^{2}} V(x) 0 \underbrace{\gamma}_{\text{parámetro}} f(x) \]

Dependiendo del valor de gamma nos modula la intensidad del potencial en algún momento. 

Tenemos: 

\[ \ddot{\psi}(x) + \left[ \gamma f(x)   + k_{0}^{2}\right] \psi(x) =0 \]


Consecuencia de la periodicidad podemos definir una expansión en \textbf{serie de Fourier}:


\[ f(x) = \sum\limits_{s=-\infty}^{\infty} C_{s} e^{-ik_{s}x} \]


Tenemos:


\[ K_{s} = \frac{2\pi s}{a}, \quad s\in \mathbb{Z} \]

Con:


\[ C_{s} = \int\limits f(x) e^{ik_{s}x}dx \]

Integramos en un periodo de la función. 


\[ C_{s} = \int\limits_{0}^{a} s f(x) e^{ik_{s}x}dx \]


Tenemos:


\[ V_{s} = - \frac{ \hbar^{2} }{2m} \gamma C_{s} \]

Vamos relacionando todo lo que tenemos hasta ahora, tenemos también nosotros lo siguiente:


Tenemos un potencial periódico. 


Del teorema de \textbf{Bloch}:


\[  \psi (x) = u(x) e^{ikx}  \]

Con:


\[  u(x) = u(x+a) \]


Como serie:

\[ u(x) = \sum\limits_{n=-\infty}^{\infty} b_{n}e^{-ik_{n}x}, \quad k_{n} = \frac{2\pi n}{a}, \quad n \in \mathbb{Z} \]




Tenemos:

\[ \psi = \sum\limits_{n=-\infty}^{\infty} b_{n}e^{ik_{n}x} e^{ikx} = e^{ikx}\sum\limits_{n=-\infty}^{\infty} b_{n}e^{-ik_{n}x} \]

Entonces ya tenemos una función de onda que prodríamos sustituirla para resolver la ecuación de segundo orden. Vamos agregar un segundo detalle. 

Nosotros sabemos cual debe de ser la solución para partícula libre:


\[ \psi(x) =b_{0}e^{ik_{0}x} \]


Donde:


\[ k_{0}^{2} \equiv \frac{2m\epsilon}{\hbar^{2}} \]


\textbf{Proponemos}: Una función de onda que tenga de manera general:

\[  \psi(x) = b_{0}e^{ikx} + \gamma \left[  e^{ikx} \sum\limits_{n\neq 0} b_{n} e^{-ik_{n}x}\right] \]

Cuando $\gamma$ es cero son queda un solución para \underline{partícula libre}.

Esta anterior la sustituimos acá:

\[ \ddot{\psi}(x) + \left[ \gamma f(x) + k_{0}^{2} \right] \psi =0 \]

Tenemos:

\begin{figure}[H]
    \centering
    \includegraphics[width=0.9\textwidth]{1.png}
    \caption{Segunda derivada de $  \psi(x) = b_{0}e^{ikx} + \gamma \left[  e^{ikx} \sum\limits_{n\neq 0} b_{n} e^{-ik_{n}x}\right]$. VL solution :)}
    \label{fig:espacio_configuracion}
\end{figure}


\[ -k^{2}b_{0}e^{ikx}+\gamma e^{ikx}\sum\limits_{n\neq 0} b_{n}e^{-ik_{n}x} (k-k_{n})^{2} + \left[ \gamma f(x) + k_{0}^{2}   \right]\left[ b_{0}e^{ikx}+\gamma e^{ikx}\sum\limits_{n\neq 0}b_{n}e^{-ik_{n}x} \right] =0\]



Tenemos:


\[ - k^{2}b_{0}+\gamma \sum\limits_{n\neq 0} b_{n}e^{-ik_{n}x} \left\{  (k-k_{n})^{2}+k_{0}^{2}  \right\} + \gamma^{2} f(x)\sum\limits_{n \neq 0} b_{n}e^{-ik_{n}x}  + \gamma f(x)b_{0}+k_{0}^{2}b_{0}=0 \]

Tenemos:


\[ b_{0} (k_{0}^{2}-k^{2}) + \gamma \left\{ f(x) b_{0} + \sum\limits_{n\neq}b_{n}e^{-ik_{n}x} \left\{ k_{0}^{2} - (k-k_{n})^{2}\right\}\right\} + \underbrace{\gamma^{2}f(x) \sum\limits_{n\neq 0} b_{n}e^{-ik_{n}x}}_{A}  =0\]



Donde:

\[ \epsilon = \frac{\hbar^{2}k_{0}^{2} }{2m} \]


Tenemos $V(x)$ débil, entonces $\gamma$ es pequeña.

Tenemos:


\[  - \frac{2m}{\hbar^{2}} V(x) = \gamma f(x)  \]


SI el potencial es intenso, entonces $\gamma$ es considerable.

Vamos a trabajar algebraicamente con un caso que si podamos resolver, que si el potencial es pequeña, entonces podemos despreciar esos ordenes.

En el examen tendríamos que resolver toda la expresión considerando la parte que ahora vamos a despreciar, es decir tendríamos que considerar también el término $\boldsymbol{A}$.

\subsection{Resolviendo}


Vamos a resolver:


\[ b_{0}(k_{0}^{2}-k^{2}) + \gamma f(x)b_{0}+ \gamma \sum\limits_{n\neq 0} \left[  k_{0}^{2} - (k_{n}-k)^{2}  \right] b_{n}e^{-ik_{n}x} =0\]

Tenemos:


\[  f(x) = \sum\limits_{s=-\infty}^{\infty} C_{s}e^{-ik_{s}x} = C_{0}e^{-ik_{s=0}x} + \sum\limits_{s \neq 0} C_{s}e^{-ik_{s}x} = C_{0} + \sum\limits_{s \neq 0} C_{s} e^{-ik_{s}x} = \sum\limits_{n \neq 0}C_{n}e^{-ik_{n}x} \]


Donde:

\[  K_{s=0} = \frac{2\pi s}{a} = 0 \]
 
Podemos elegir por conveniencia que el potencial valga \textbf{cero}.  

Entonces: 

\[ \cancelto{0}{C_{0}} \]

\begin{figure}[H]
    \centering
    \includegraphics[width=0.9\textwidth]{2.png}
    \caption{.}
    \label{fig:espacio_configuracion}
\end{figure}

Tenemos:



\[ \Rightarrow \quad b_{0}( k_{0}^{2} - k^{2}) + \gamma \left\{ \sum\limits_{n\neq 0} b_{0}C_{n}e^{-ik_{n}x} + \sum\limits_{n\neq 0} b_{n}e^{-ik_{n}x} \left[  k_{0}^{2}-(k_{n}-k)^{2}  \right]   \right\}  \]


Multiplicamos todo por:


\[ e^{ik_{m}x}, \quad m \in \mathbb{Z}\, \text{arbitrario} \]

e integrar en una celda.

Escencialmente los términos $C_{0}$ de la serie están relacionados con el potencial.


Tenemos:


\[  \int\limits_{0}^{a} (k_{0}^{2} - k^{2}) e^{ik_{m}x} dx + \gamma \int\limits_{0}^{a} \sum\limits_{n\neq 0} b_{0}c_{n}e^{-ik_{n}x} e^{ik_{m}x} dx + \gamma \int\limits_{0}^{a} \sum\limits_{n\neq 0} b_{n} \left[  k_{0}^{2} - (k_{n}-k)^{2}  \right]e^{-ik_{n}x}e^{ik_{m}x} dx=0   \]



Tenemos:

\[ e^{i k_{m}x} = e^{i  \frac{2\pi mx}{a} } \]


\[  \int\limits_{0}^{a} b_{0} (k_{0}^{2} - k^{2})e^{ik_{m}x}dx  \quad \to \quad \int\limits_{0}^{a} e^{ik_{m}x} dx = \left. \left( \frac{1}{ik_{m}} \right) e^{ik_{m}x} \right|_{0}^{a}  =\left.  \frac{ a  }{i 2\pi m} e^{ik_{m}x}\right|_{0}^{a}  \]


\[ k_{m} = \frac{2\pi }{a} m, \quad m \in \mathbb{Z}, \quad \left\{ \begin{array}{cc}
0 & m \neq 0\\
a & m=0
\end{array} \right. \]

Es decir:


\[ \left(  \frac{a}{ i2\pi m }  \right) \left[  e^{ i\frac{2\pi m}{a} a} - \cancelto{1}{ e^{i \frac{2\pi m}{a}(0)  }   }  \right]  \]


\[ = \left(   \frac{a}{i 2\pi m  } \right)(0) \]


Tenemos:


\[ \int\limits_{0}^{a} e^{i_{m}x} dx = \int\limits_{0}^{a} dx = a   \]


Tenemos:


\[  \int\limits_{0}^{a} e^{i (k_{m}-k_{n})x }dx = \left. \frac{1}{ i (k_{m}-k_{n})   } e^{i (k_{m}-k_{n})  x} \right|_{0}^{a}  = \frac{  e^{ i (k_{m} - k_{n})a } -1 }{ i ( k_{m} - k_{n} )  } =   \left\{ \begin{array}{cc}
0 & m \neq n\\
a & m=n
\end{array} \right.  \]

\[ k_{m} - k_{n} = \frac{ 2\pi (m-n)   }{a}  \]

\begin{figure}[H]
    \centering
    \includegraphics[width=0.9\textwidth]{3.png}
    \caption{.}
    \label{fig:espacio_configuracion}
\end{figure}

\begin{figure}[H]
    \centering
    \includegraphics[width=0.9\textwidth]{4.png}
    \caption{.}
    \label{fig:espacio_configuracion}
\end{figure}









\end{document}
